\section{Cambio en la estructura del mercado: de un centro basado en apps a un centro basado en la interfaz del agente}

Uno de los juicios más polémicos de la entrevista es que «el 80\% de las apps podría desaparecer o debilitarse». Dicho con más precisión: «muchas apps de interacción ligera serán absorbidas por la capa de orquestación del agente».

\subsection{Por qué la «desaparición de las apps» es, en esencia, una migración de la capa de interacción}

Cuando el agente puede operar directamente la API o la GUI, el usuario ya no necesita entrar repetidamente a una app de función única. El centro de valor se desplazará de «la posesión de la interfaz» hacia «la posibilidad de invocar la capacidad».

\begin{importantbox}{Valor del producto que se revalúa}
\begin{itemize}
    \item el valor de la UI puntual disminuye;
    \item el valor de la disponibilidad de datos y de la confiabilidad de la API aumenta;
    \item la orquestación de permisos y los mecanismos de liquidación se convierten en la nueva infraestructura;
    \item los servicios agent-friendly obtendrán la ventaja de ser los primeros.
\end{itemize}
\end{importantbox}

\begin{knowledgebox}{Dos vías de transformación para las empresas}
\begin{enumerate}
    \item convertir la app existente en un producto de API estable;
    \item aceptar al agente como primer punto de entrada y rediseñar la estrategia de cobro y control de riesgo.
\end{enumerate}
Si no se hace ninguna de las dos cosas, la empresa será devorada poco a poco por una «capa sustituta programable».
\end{knowledgebox}

\begin{warningbox}{Conflicto entre las estrategias antirrastreo y el valor para el usuario}
La entrevista menciona el fuerte bloqueo que las plataformas imponen a los bots. A corto plazo esto puede prevenir el abuso, pero si bloquea también el acceso legítimo de los agentes de usuarios reales, a largo plazo puede empujar a los usuarios hacia plataformas alternativas más abiertas.
\end{warningbox}

\begin{figure}[H]
\centering
\includegraphics[width=0.9\textwidth]{frames/frame-023622.jpg}
\caption{Discusión sobre la forma del producto: Heartbeat, activación proactiva y «el servicio como interfaz»\protect\footnotemark}
\end{figure}
\footnotetext{Intervalo de tiempo en el video: 02:36:15--02:36:35.}

\subsection{Resumen de la sección}
La esencia de la controversia sobre OpenClaw no es «si debe haber apps o no», sino «quién posee el punto de entrada de la tarea». En cuanto el agente se convierta en el punto de entrada predeterminado, el activo central de las empresas de software migrará de la interfaz hacia la interfaz de programación y los datos.

